% GENERATED FILE. Edit canonical lessons, then run npm run generate:books.
% canonical-source-hash: fnv1a64:2863dcf13529bf1c
% canonical-lessons: ML-C354-prasamgam, ML-C354-madhyamam, ML-C354-prasamgikkuka, ML-C354-kayyatikkuka, ML-C354-paramarsikkuka

\chapter{A Speech and Its Applause}
\label{ch:ml-a-speech-and-its-applause}
\hlchaptermodality{354}

\begin{chapteropening}
\textbf{By the end of this chapter:} \emph{I can say a speech, a medium, to give a speech, to applaud and to mention.}

Complete the last lesson of chapter 354: I can say a speech, a medium, to give a speech, to applaud and to mention.
\end{chapteropening}

\section[\texorpdfstring{\ml{പ്രസംഗം} (prasaṁgaṁ)}{പ്രസംഗം (prasaṁgaṁ)}]{\ml{പ്രസംഗം} (prasaṁgaṁ) — a speech}

\label{lesson:ML-C354-prasamgam}



\begin{quote}
Before the new one: say the Malayalam for to circulate, then the Malayalam for to spread around.
\end{quote}

\subsection*{You'll want to know: \ml{പ്രസംഗം}}
\textbf{\ml{പ്രസംഗം}} — \emph{prasaṁgaṁ} — \textquotedblleft{}a speech\textquotedblright{}.

\begin{quote}
\ml{പ്രിൻസിപ്പലിന്റെ} \ml{പ്രസംഗം} \ml{വളരെ} \ml{നീണ്ടതായിരുന്നു}.
\end{quote}

\begin{quote}
\emph{prinsippalinṟe prasaṁgaṁ vaḷare nīṇṭatāyirunnu.}
\end{quote}

\begin{quote}
The principal's speech was very long.
\end{quote}

A speech is given: \ml{പ്രസംഗം} \ml{നടത്തുക}.

\begin{cousinweb}
\textbf{\ml{പ്രസംഗം}} (\emph{prasaṁgaṁ}) is borrowed from Sanskrit \emph{prasaṅga}, \textquotedblleft{}adherence, attachment, connection\textquotedblright{} (Monier-Williams). Malayalam narrowed it to a talk given to an audience.
\end{cousinweb}

\begin{grammarlens}[title={What you've built}]
One more word for telling what happened.
\end{grammarlens}

\subsection*{Your turn}
\begin{itemize}
\raggedright
  \item \emph{Say it:} \emph{prasaṁgaṁ}
  \item \emph{Say it:} \emph{prasaṁgaṁ}, once more
  \item \emph{Say it:} \emph{prinsippalinṟe prasaṁgaṁ vaḷare nīṇṭatāyirunnu.} — The principal's speech was very long.
  \item \emph{Recall:} say the Malayalam for to circulate, then the Malayalam for to spread around, then say \emph{prasaṁgaṁ} again
\end{itemize}

\begin{tcolorbox}[breakable,colback=teal!4,colframe=teal!35!black,title={Before you move on}]
What does \ml{പ്രസംഗം} mean? (\textquotedblleft{}A speech\textquotedblright{}.) Say it once more.
\end{tcolorbox}
\section[\texorpdfstring{\ml{മാധ്യമം} (mādhyamaṁ)}{മാധ്യമം (mādhyamaṁ)}]{\ml{മാധ്യമം} (mādhyamaṁ) — a medium; the media}

\label{lesson:ML-C354-madhyamam}



\begin{quote}
Before the new one: say the Malayalam for to spread around, then the Malayalam for a speech.
\end{quote}

\subsection*{You'll want to know: \ml{മാധ്യമം}}
\textbf{\ml{മാധ്യമം}} — \emph{mādhyamaṁ} — \textquotedblleft{}a medium; the media\textquotedblright{}.

\begin{quote}
\ml{മാധ്യമങ്ങൾ} \ml{ആ} \ml{സംഭവം} \ml{വലിയ} \ml{വാർത്തയാക്കി}.
\end{quote}

\begin{quote}
\emph{mādhyamaṅṅaḷ ā sambhavaṁ valiya vārttayākki.}
\end{quote}

\begin{quote}
The media made that incident big news.
\end{quote}

Plural: \ml{മാധ്യമങ്ങൾ}, \emph{mādhyamaṅṅaḷ}.

In the plural, \ml{മാധ്യമങ്ങൾ}, it means the press and television together.

\begin{cousinweb}
\textbf{\ml{മാധ്യമം}} (\emph{mādhyamaṁ}) is borrowed from Sanskrit \emph{mādhyama}, \textquotedblleft{}relating to the middle, central\textquotedblright{} (Monier-Williams), from \emph{madhya}, \textquotedblleft{}middle\textquotedblright{}. A medium stands in the middle and carries news between people.
\end{cousinweb}

\begin{grammarlens}[title={What you've built}]
One more word for telling what happened.
\end{grammarlens}

\subsection*{Your turn}
\begin{itemize}
\raggedright
  \item \emph{Say it:} \emph{mādhyamaṁ}
  \item \emph{Say it:} \emph{mādhyamaṁ}, once more
  \item \emph{Say it:} \emph{mādhyamaṅṅaḷ ā sambhavaṁ valiya vārttayākki.} — The media made that incident big news.
  \item \emph{Recall:} say the Malayalam for to spread around, then the Malayalam for a speech, then say \emph{mādhyamaṁ} again
\end{itemize}

\begin{tcolorbox}[breakable,colback=teal!4,colframe=teal!35!black,title={Before you move on}]
What does \ml{മാധ്യമം} mean? (\textquotedblleft{}A medium; the media\textquotedblright{}.) Say it once more.
\end{tcolorbox}
\section[\texorpdfstring{\ml{പ്രസംഗിക്കുക} (prasaṁgikkuka)}{പ്രസംഗിക്കുക (prasaṁgikkuka)}]{\ml{പ്രസംഗിക്കുക} (prasaṁgikkuka) — to give a speech}

\label{lesson:ML-C354-prasamgikkuka}



\begin{quote}
Before the new one: say the Malayalam for a speech, then the Malayalam for a medium.
\end{quote}

\subsection*{You'll want to know: \ml{പ്രസംഗിക്കുക}}
\textbf{\ml{പ്രസംഗിക്കുക}} — \emph{prasaṁgikkuka} — \textquotedblleft{}to give a speech\textquotedblright{}.

\begin{quote}
\ml{വാർഷികത്തിന്} \ml{പ്രിൻസിപ്പൽ} \ml{പ്രസംഗിച്ചു}.
\end{quote}

\begin{quote}
\emph{vārṣikattinŭ prinsippal prasaṁgiccu.}
\end{quote}

\begin{quote}
The principal gave a speech at the anniversary.
\end{quote}

Past: \ml{പ്രസംഗിച്ചു}, \emph{prasaṁgiccu}.

Where takes -\ml{ഇൽ}: \ml{യോഗത്തിൽ} \ml{പ്രസംഗിച്ചു}, spoke at the meeting.

\begin{cousinweb}
\textbf{\ml{പ്രസംഗിക്കുക}} (\emph{prasaṁgikkuka}) is built on \ml{പ്രസംഗം} (\emph{prasaṁgaṁ}), \textquotedblleft{}a speech\textquotedblright{}, from Sanskrit \emph{prasaṅga} (Monier-Williams), with the Malayalam verb ending -\ml{ിക്കുക}, the pattern that makes \ml{ചിന്തിക്കുക} from \ml{ചിന്ത}.
\end{cousinweb}

\begin{grammarlens}[title={What you've built}]
One more word for telling what happened.
\end{grammarlens}

\subsection*{Your turn}
\begin{itemize}
\raggedright
  \item \emph{Say it:} \emph{prasaṁgikkuka}
  \item \emph{Say it:} \emph{prasaṁgikkuka}, once more
  \item \emph{Say it:} \emph{vārṣikattinŭ prinsippal prasaṁgiccu.} — The principal gave a speech at the anniversary.
  \item \emph{Recall:} say the Malayalam for a speech, then the Malayalam for a medium, then say \emph{prasaṁgikkuka} again
\end{itemize}

\begin{tcolorbox}[breakable,colback=teal!4,colframe=teal!35!black,title={Before you move on}]
What does \ml{പ്രസംഗിക്കുക} mean? (\textquotedblleft{}To give a speech\textquotedblright{}.) Say it once more.
\end{tcolorbox}
\section[\texorpdfstring{\ml{കയ്യടിക്കുക} (kayyaṭikkuka)}{കയ്യടിക്കുക (kayyaṭikkuka)}]{\ml{കയ്യടിക്കുക} (kayyaṭikkuka) — to applaud}

\label{lesson:ML-C354-kayyatikkuka}



\begin{quote}
Before the new one: say the Malayalam for a medium, then the Malayalam for to give a speech.
\end{quote}

\subsection*{You'll want to know: \ml{കയ്യടിക്കുക}}
\textbf{\ml{കയ്യടിക്കുക}} — \emph{kayyaṭikkuka} — \textquotedblleft{}to applaud\textquotedblright{}.

\begin{quote}
\ml{പ്രസംഗം} \ml{കഴിഞ്ഞപ്പോൾ} \ml{എല്ലാവരും} \ml{കയ്യടിച്ചു}.
\end{quote}

\begin{quote}
\emph{prasaṁgaṁ kaḻiññappōḷ ellāvaruṁ kayyaṭiccu.}
\end{quote}

\begin{quote}
When the speech was over, everyone applauded.
\end{quote}

Past: \ml{കയ്യടിച്ചു}, \emph{kayyaṭiccu}.

Literally to strike the hands: \ml{കൈ}, hand, + \ml{അടിക്കുക}, to strike.

\begin{cousinweb}
\textbf{\ml{കയ്യടിക്കുക}} (\emph{kayyaṭikkuka}) is a native Dravidian compound of \ml{കൈ} (\emph{kai}), \textquotedblleft{}hand\textquotedblright{} (DEDR 2023), and \ml{അടിക്കുക} (\emph{aṭikkuka}), \textquotedblleft{}to strike, to beat\textquotedblright{} (DEDR 77): to applaud is to strike the hands together.
\end{cousinweb}

\begin{grammarlens}[title={What you've built}]
One more word for telling what happened.
\end{grammarlens}

\subsection*{Your turn}
\begin{itemize}
\raggedright
  \item \emph{Say it:} \emph{kayyaṭikkuka}
  \item \emph{Say it:} \emph{kayyaṭikkuka}, once more
  \item \emph{Say it:} \emph{prasaṁgaṁ kaḻiññappōḷ ellāvaruṁ kayyaṭiccu.} — When the speech was over, everyone applauded.
  \item \emph{Recall:} say the Malayalam for a medium, then the Malayalam for to give a speech, then say \emph{kayyaṭikkuka} again
\end{itemize}

\begin{tcolorbox}[breakable,colback=teal!4,colframe=teal!35!black,title={Before you move on}]
What does \ml{കയ്യടിക്കുക} mean? (\textquotedblleft{}To applaud\textquotedblright{}.) Say it once more.
\end{tcolorbox}
\section[\texorpdfstring{\ml{പരാമർശിക്കുക} (parāmarśikkuka)}{പരാമർശിക്കുക (parāmarśikkuka)}]{\ml{പരാമർശിക്കുക} (parāmarśikkuka) — to mention; to refer to}

\label{lesson:ML-C354-paramarsikkuka}



\begin{quote}
Before the new one: say the Malayalam for to give a speech, then the Malayalam for to applaud.
\end{quote}

\subsection*{You'll want to know: \ml{പരാമർശിക്കുക}}
\textbf{\ml{പരാമർശിക്കുക}} — \emph{parāmarśikkuka} — \textquotedblleft{}to mention; to refer to\textquotedblright{}.

\begin{quote}
\ml{പ്രസംഗത്തിൽ} \ml{അവൾ} \ml{അമ്മയെ} \ml{പരാമർശിച്ചു}.
\end{quote}

\begin{quote}
\emph{prasaṁgattil avaḷ ammaye parāmarśiccu.}
\end{quote}

\begin{quote}
In the speech she mentioned her mother.
\end{quote}

Past: \ml{പരാമർശിച്ചു}, \emph{parāmarśiccu}.

The noun is \ml{പരാമർശം}. What is mentioned takes the object ending: \ml{അപകടത്തെ} \ml{പരാമർശിച്ചു}, mentioned the accident.

\begin{cousinweb}
\textbf{\ml{പരാമർശിക്കുക}} (\emph{parāmarśikkuka}) is built from Sanskrit \emph{parāmarśa}, \textquotedblleft{}seizing, taking hold of\textquotedblright{} (Monier-Williams), with the Malayalam verb ending -\ml{ിക്കുക}: to mention a thing is to take hold of it in words.
\end{cousinweb}

\begin{grammarlens}[title={What you've built}]
One more word for telling what happened.
\end{grammarlens}

\subsection*{Your turn}
\begin{itemize}
\raggedright
  \item \emph{Say it:} \emph{parāmarśikkuka}
  \item \emph{Say it:} \emph{parāmarśikkuka}, once more
  \item \emph{Say it:} \emph{prasaṁgattil avaḷ ammaye parāmarśiccu.} — In the speech she mentioned her mother.
  \item \emph{Recall:} say the Malayalam for to give a speech, then the Malayalam for to applaud, then say \emph{parāmarśikkuka} again
\end{itemize}

\begin{tcolorbox}[breakable,colback=teal!4,colframe=teal!35!black,title={Before you move on}]
What does \ml{പരാമർശിക്കുക} mean? (\textquotedblleft{}To mention; to refer to\textquotedblright{}.) Say it once more.
\end{tcolorbox}
